%\chapter*{\nomname} % Заголовок
%\addcontentsline{toc}{chapter}{\nomname}  % Добавляем его в оглавление
\noskipnonumchapter{\nomname}
\noindent
%\begin{longtabu} to \dimexpr \textwidth-5\tabcolsep {r X}
\begin{longtabu} to \textwidth {r X}

% Аэродинамические параметры
\(M\) & число Маха \\
\(Re\) & число Рейнольдса \\
\(C_x\) & коэффициент сопротивления \\
\(C_y\) & коэффициент подъёмной силы \\
\(\alpha\) & угол атаки \\
\(\beta\) & угол скольжения \\

% Аббревиатуры
ЛА & летательный аппарат \\
ТРДД & турбореактивный двухконтурный двигатель \\
САХ & средняя аэродинамическая хорда \\
CFD & computational fluid dynamics, вычислительная гидродинамика \\
RANS & Reynolds-averaged Navier–Stokes, осреднённые по Рейнольдсу уравнения Навье–Стокса \\
LES & large eddy simulation, метод моделирования крупных вихрей \\
UAV & unmanned aerial vehicle, беспилотный летательный аппарат \\

\end{longtabu}
\addtocounter{table}{-1}% Нужно откатить на единицу счетчик номеров таблиц, так как предыдующая таблица сделана для удобства представления информации по ГОСТ
